% ============================================================================
%  CHƯƠNG 5. KẾT LUẬN, HẠN CHẾ VÀ HƯỚNG PHÁT TRIỂN
% ============================================================================
\clearpage
\section{Kết luận, hạn chế và hướng phát triển}
\label{sec:conclusion}

\subsection{Kết luận}

\textcolor{badred}{[Điền sau E1--E5: trả lời Q1 (LoRA có đạt sens@spec0,90 $\ge 0{,}60$ trên VinBigData test
không) và Q2 (cải thiện có còn trên JSRT và khi che nốt không), theo cả hai hướng có thể xảy ra.]}

Dù kết quả theo hướng nào, đồ án sẽ đóng góp một benchmark có kiểm soát cho VLM trên bài toán nốt phổi, một
phân tích định lượng về chế độ báo bình thường khi tinh chỉnh với lớp dương hiếm, và một quy trình kiểm toán
``nhìn đúng chỗ'' dùng được cho các trợ lý đọc ảnh khác.

\subsection{Hạn chế}

\begin{itemize}
  \item \textbf{Nhãn có nhiễu:} nhãn NIH rút từ báo cáo; nhãn VinBigData có bất đồng giữa các bác sĩ.
  \item \textbf{Cỡ mẫu nhỏ:} khoảng 82 ảnh dương ở tập test VinBigData, nên khoảng tin cậy rộng.
  \item \textbf{Bootstrap theo ảnh, không gộp theo bệnh nhân:} nếu một bệnh nhân có nhiều ảnh, khoảng tin cậy có
        thể hẹp hơn thực tế.
  \item \textbf{Chia VinBigData theo ảnh:} bản Kaggle không có mã bệnh nhân, nên ảnh của cùng một bệnh nhân (nếu
        có) có thể nằm ở cả train và test, làm kết quả trong phân phối lạc quan hơn.
  \item \textbf{Thiếu kích thước thật:} VinBigData không có kích thước nốt tính bằng mm; chỉ JSRT có.
  \item \textbf{Ngoài phân phối mạnh:} JSRT là phim scan cũ.
  \item \textbf{Mô hình cỡ nhỏ}, bộ mã hoá ảnh đông cứng; kết luận không tự động mở rộng sang VLM lớn.
  \item \textbf{Không có mô bệnh học:} trợ lý phát hiện nốt/khối, không phân biệt lành/ác (trừ phân tích phụ trên JSRT).
  \item \textbf{Rò rỉ dữ liệu:} MIMIC-CXR, VinDr-CXR và NIH được tái dùng rộng rãi để huấn luyện
        VLM~\cite{johnson2019mimic,nguyen2025vindrvqa}, nên kết quả zero-shot trên các tập này có thể bị thổi phồng.
        JSRT ít bị dùng hơn nhưng không đảm bảo hoàn toàn.
\end{itemize}

\subsection{An toàn và đạo đức}

Một trợ lý yếu nhưng trả lời tự tin có thể kéo người đọc theo hướng sai~\cite{bernstein2023incorrect,lee2023humanai}.
Vì vậy kiểm toán ``nhìn đúng chỗ'' có ý nghĩa an toàn, không chỉ học thuật. \textbf{Trợ lý trong đồ án là nguyên
mẫu nghiên cứu, không dùng cho quyết định lâm sàng} và không được trình bày như công cụ sẵn sàng triển khai. Đồ
án chỉ dùng dữ liệu công khai đã ẩn danh và không phân phối lại ảnh.

\subsection{Hướng phát triển}

\begin{itemize}
  \item \textbf{VLM có hỗ trợ vùng (C5):} detector đề xuất vùng, VLM xác nhận và giải thích, theo kiểu agent
        MedRAX~\cite{fallahpour2025medrax}.
  \item \textbf{Từ phát hiện đến báo cáo và theo dõi:} sinh câu mô tả nốt (vị trí, kích thước, tính chất) và gợi ý
        chụp CT theo guideline Fleischner, có RAG; các LLM hiện vẫn áp guideline chưa đáng tin cậy~\cite{gamble2024fleischner}.
  \item \textbf{Nguy cơ ung thư đa phương thức:} kết hợp CXR với dữ liệu lâm sàng, so với CXR-LC~\cite{lu2020cxrlc}.
  \item \textbf{Câu hỏi mở:} bộ mã hoá ảnh đông cứng có phải nút thắt không? Cần thí nghiệm mở khoá một phần bộ mã
        hoá, hoặc probe patch-local trên mô hình sau LoRA.
\end{itemize}
